%LTeX: language=fr-FR
\documentclass{article}
\usepackage{tasks}

\usepackage{amsmath,amssymb,amsfonts}
\usepackage{graphicx}
\usepackage{amsthm}
\usepackage[french]{babel}
\usepackage{enumitem}
\usepackage{currfile}

\usepackage{tikz}
\usetikzlibrary{calc,intersections}
\usepackage{qrcode}
\usepackage{mathrsfs}
\usepackage{enumitem}
\usepackage[french]{babel}
\usepackage{tcolorbox}
\usepackage{fancyhdr}
\usepackage{multicol}
\usepackage[headsep=3mm,footskip=5mm,includefoot,includehead]{geometry}
%showframe
\usepackage[lastpage,user]{zref}

\newtheorem*{theorem*}{Théorème}

\setlength{\parindent}{0cm}

\newenvironment{clist}{
  \begin{center}
  \begin{inparaenum}}{
    \end{inparaenum}
\end{center}}

\pagestyle{empty}
\graphicspath{ {Images/} }
\input{\currfiledir ../def.tex}
\def \Chapitre {Chapitre 3 -- Fonctions affines}
\def \ChapitreCourt {Ch. 3 -- Fonctions affines}


\geometry{left=1.5cm, right=1.5cm, top=1cm, bottom=1cm}
\setlength{\headheight}{14.0pt}

\theoremstyle{definition}
\newtheorem{exercice}{Exercice}

\begin{document}
\pagestyle{fancy}

\fancyhf{} % clear existing header/footer entries
% We don't need to specify the O coordinate
\fancyhead[L]{\large \textbf{\Niveau}}
\fancyhead[C]{\large \textbf{\Chapitre}}

\large

\begin{paragraph}{Pré-requis}
  \begin{itemize}
    \item Généralités sur les fonctions (chapitre 1) : image, antécédent, courbe, résolution graphique de $f(x)=k$ et $f(x)<k$, tableau de variations
    \item Équations du premier degré (questions flash)
    \item Calcul avec les relatifs et les fractions
    \item Lire et placer des points dans un repère
  \end{itemize}
\end{paragraph}

\begin{paragraph}{Contenus}
  \begin{itemize}
    \item Fonction affine $x \mapsto ax+b$ ; cas particuliers : fonction linéaire ($b=0$), fonction constante ($a=0$)
    \item Représentation graphique : une droite d'équation réduite $y=ax+b$
    \item Coefficient directeur $a$ et ordonnée à l'origine $b$
    \item Calcul du coefficient directeur à partir de deux points
    \item Sens de variation selon le signe de $a$
    \item Signe de $ax+b$ et tableau de signes
  \end{itemize}
\end{paragraph}

\begin{paragraph}{Objectifs}
  \begin{itemize}
    \item Reconnaître une fonction affine et identifier $a$ et $b$,
    \item Modéliser une situation par une fonction affine (tarifs, coûts, évolutions régulières),
    \item Tracer une droite donnée par son équation réduite, ou par un point et son coefficient directeur,
    \item Lire graphiquement l'équation réduite d'une droite,
    \item Calculer le coefficient directeur à partir des coordonnées de deux points,
    \item Interpréter $a$ comme l'augmentation de $f(x)$ quand $x$ augmente de $1$,
    \item Relier signe de $a$, sens de variation et allure de la droite,
    \item Déterminer le signe de $ax+b$ et dresser son tableau de signes,
    \item Résoudre $f(x)=k$ et $f(x)<k$ par le calcul et graphiquement ; comparer deux tarifs (en exercice).
  \end{itemize}
\end{paragraph}

\end{document}
